% ===== 原文 pp.245–261（PDF p.260–276）=====
% 第三篇扉页（p.245）、第 6 章标题与 §6.1（pp.247–254）、§6.2 开头（pp.255–261）
% tag 区间：6.1–6.2（(6.1) 在 p.249；(6.2) 在 p.252）
% 本块写篇标题与章标题；§6.1、§6.2 及其 6.1.1–6.1.3、6.2.1 小节标题由本块写出。
% 图：第 6 章 Figure 1(p.255)、2(p.259)、3(p.260)、4(p.261) → work/figures/f6_1..f6_4.png。
% 块末 Lemma 6.9 的证明在 p.262 续（下一块 c06b 开头 "To show this independence …"），
% 故本块把该证明写成 \para{证明。} 起的普通段落、不闭合 proof 环境，由 c06b 在末尾补 □。
% [原文 p.245]
\part{第三篇\quad Gross--Siebert 纲领}

% [原文 p.246]（原文此页为空白页）

\section{Gross--Siebert 纲领与二维结果}

% [原文 p.247]
我们已经看到，至少对 $\PP^2$ 而言，热带几何为镜像对称的 A-模型与 B-模型给出了一个非常令人满意的解释。本章我们将探讨如何用类似的想法来解释更复杂的簇（尤其是 Calabi–Yau 流形）的镜像对称。

这里描述的想法是与 Bernd Siebert 的合作工作，我们于 2001 年开始了这一纲领（见 [47]、[48]、[40]、[51]、[49]）。目前已有若干解释该纲领不同方面的综述文章：特别参见 [41]，它探讨了该纲领如何从对 Strominger–Yau–Zaslow（SYZ）猜想的研究中生长出来；以及 [3] 的第 6、7 章，那里对 SYZ 猜想的若干方面有更详尽的描述。

这里我们采取更为非历史的观点，通过与本书已涵盖的想法作类比来叙述该纲领。我们先大致解释第 1 章引入的热带流形与具有有效反典范类的"良好"退化族之间的对应——这包括镜像对称最著名的例子 Calabi–Yau 情形。这就把我们引向从热带流形重构这样一族簇的基本问题。

本章接下来两节将说明我与 Siebert 在 [49] 中对最简单的一维情形（即二维情形）给出的该问题的解答。除依赖 [45] 的结果外，论证是完整的。我这里给出的证明更多是 Kontsevich–Soibelman [70] 的途径（它只处理二维情形）与 [49] 的途径（它处理所有维数，但显著更难）的混合体。

% [Tutor Note: 章标题原文为 "The program and two-dimensional results"，本译按 BRIEF 作 "Gross--Siebert 纲领与二维结果"。本章 §6.1 是纲领总览，§6.2–6.3 只处理 $\dim B=2$；本块（pp.245–261）覆盖 §6.1 全部与 §6.2 开头，其余由 c06b 续。]

\subsection{纲领}

第 4 章的主要内容是顶点位于 $M_{\mathbb{Q}}$ 的热带曲线与定义在域 $k((t))$ 上的代数曲线之间的对应定理。我们在那里看到，某种强整性使得热带曲线对应的代数曲线定义在 $k((t))$ 上，而不仅仅是定义在 $\overline{k((t))}$ 上。

一般地，人们未必愿意处理顶点落在 $M_{\mathbb{Q}}$ 中的热带曲线；此时应当在域
\[
  K=\Bigl\{\sum_i a_it^{r_i}\ \Big|\ a_i\in k,\ r_i\in\mathbb{R}\ \text{构成满足}\ r_i\to\infty\ \text{的序列}\Bigr\}
\]
上工作。这个域带有赋值 $\nu:K^{\times}\to\mathbb{R}$，由
\[
  \nu\Bigl(\sum_i a_it^{r_i}\Bigr)=\min\{r_i\}
\]
给出。

% [原文 p.248]
给定簇 $X\subseteq(K^{\times})^n$，我们便可定义 $X$ 的热带化 $\operatorname{Trop}(X)$ 为
\[
  \operatorname{Trop}(X)=\{(\nu(x_1),\dots,\nu(x_n))\mid(x_1,\dots,x_n)\in X\}\subseteq\mathbb{R}^n.
\]
可以证明 $\operatorname{Trop}(X)$ 是一种热带簇。不过我们不打算精确定义这里"热带簇"的含义；关于热带簇的定义，文献中有很多讨论，例如见 [96]、[81]。当然，热带簇包含本书已经研究过的热带曲线。

% [Tutor Note: 域 $K$ 的赋值 $\nu(\sum a_it^{r_i})=\min\{r_i\}$ 说明顶点在 $M_{\mathbb{Q}}$ 的热带曲线自然只给出 $\nu$ 取有理值的曲线；若顶点落在 $M_{\mathbb{Z}}$（即第 4 章所说的"强整性"），则对应曲线可定义在较小的 $k((t))$ 上，这正是上一节末句"而不仅仅是 $\overline{k((t))}$"的含义。]

于是人们自然要问：$\mathbb{R}^n$ 中的热带簇何时能实现为 $(K^{\times})^n$ 中某个簇的热带化？第 4 章本质上回答了 $\mathbb{R}^2$ 中亏格零热带曲线的这一问题；事实上，由 Mikhalkin [80] 的工作，同样的结论对 $\mathbb{R}^2$ 中任意亏格的曲线也成立。然而一般地，在 $\mathbb{R}^n$（$n\ge3$）中，只有亏格零的热带曲线才来自真正的曲线；更高亏格时未必如此：这一方向的一些结果见 [106]。对于更高维的热带簇，人们知之甚少。

让我们稍作观念上的转变。上一段中 $\mathbb{R}^n$ 是外围空间，是 $(K^{\times})^n$ 的"热带化"，我们关心的是 $\mathbb{R}^n$ 的哪些热带子簇是 $(K^{\times})^n$ 中簇的热带化。现在我们反过来，尝试改变外围空间。或许存在更有趣的选择，对应于更有趣的簇。特别地，若把 $\mathbb{R}^n$ 换成定义 1.22 意义下的热带仿射流形 $B$，我们就得到更大的外围簇集合。如第 1 章所见，热带仿射流形正是仍然可以谈论热带曲线的流形。麻烦在于紧热带仿射流形没有多少有趣的例子（$\mathbb{R}^n/\Gamma$，$\Gamma$ 为格，是这样一例，对应于复环面），因此我们需要允许带奇点的热带仿射流形，更确切地说，即定义 1.27 中所谓的\emph{热带流形}。

于是设 $B$ 是热带流形。此时自然会尝试给这样的流形配上一个定义在域 $K$ 上的簇。这个簇与热带对象之间究竟有何联系，并不一目了然；我们实际上将回避这一问题，因为要给出恰当的回答需要在刚性解析空间范畴中工作（见 [70]）。不过，若进一步假定 $B$ 是整的，我们就能给出一个动机合理的猜测。使用整热带流形的好处在于：我们可以期望相应的簇定义在子域 $k((t))$ 上，正如第 4 章的情形那样。

借用第 4 章的提示，我们实际上应当尝试给 $B$ 配上定义在 $\operatorname{Spec} k[t]$ 上的概形 $X$。一般纤维 $X_{\eta}$ 将是所要的 $k((t))$ 上的簇，而闭点上的纤维 $X_0$ 则是一个退化簇。这应当与第 4 章的情形气味相投——那里我们考虑了 $\PP^2$ 或其他环面曲面的退化。那些退化来自 $\mathbb{R}^2$ 的多面体分解。所以我们应当预期 $X_0$ 的构造涉及某种选择。由类比可知，我们合理地需要用到定义 1.27 中 $B$ 的多面体分解 $\mathcal{P}$。这就是 Gross–Siebert 纲领的基本背景。

在本节余下部分，我们将概述 Gross–Siebert 纲领的一些基本想法，但不给出任何技术细节。

% [原文 p.249]
\subsubsection{扇形图像，或 A-模型}

于是，给定 $(B,\mathcal{P})$，我们先来猜测我们希望给 $(B,\mathcal{P})$ 配上的退化 $X\to\operatorname{Spec} k[t]$ 的中心纤维 $X_0$ 应具有的形式。再次以第 4 章为指引，回忆我们考虑过对 $(M_{\mathbb{R}},\mathcal{P})$。$\mathcal{P}$ 的每个顶点 $v$ 定义了一个扇 $\Sigma_v$，它定义环面簇 $D_v$，而 $X_0$ 是一族不可约分支之并 $X_0=\bigcup_v D_v$，其中 $v$ 遍历 $\mathcal{P}$ 的所有顶点。$X_0$ 的各不可约分支按照 $\mathcal{P}$ 的组合学所决定的方式粘合起来。

要点在于，从整热带流形 $(B,\mathcal{P})$ 我们得到完全相同的结构。对每个胞腔 $\tau\in\mathcal{P}$，我们得到由 $B$ 上的扇结构给出的扇 $\Sigma_\tau$。对应于 $\Sigma_\tau$ 的是一个环面簇 $X_\tau$（我们现在用 $X_\tau$ 代替第 4 章中的 $D_\tau$）。特别地，对每个顶点 $v\in\mathcal{P}$，我们有环面簇 $X_v$。此外，若 $v,w$ 是 $\tau\in\mathcal{P}$ 的顶点，则 $\Sigma_\tau$ 可以描述为 $\Sigma_v$ 与 $\Sigma_w$ 的商扇，于是 $X_\tau$ 自然地是 $X_v$ 与 $X_w$ 的环面层。若 $\tau$ 是 $\mathcal{P}$ 中包含 $v$ 与 $w$ 的最小胞腔，则我们沿 $X_\tau\subseteq X_v$ 与 $X_\tau\subseteq X_w$ 这两个副本把 $X_v$ 与 $X_w$ 粘合起来。

这样我们就构造出簇 $X_0=X_0(B,\mathcal{P})$，它是一族环面簇之并。例如，把它应用于例 1.28 (3)（其中 $(B,\mathcal{P})$ 由 $\partial\Xi_3$ 给出），便得到八个 $\PP^2$ 之并，粘合成一个八面体。

一般地，情形实际上更复杂一些：$X_\tau\subseteq X_v$ 与 $X_\tau\subseteq X_w$ 的保环面层的等同方式有选择的余地。这给出可能的粘合的某种模，但总存在一个典范选择，即等同两个 $X_\tau$ 副本中大环面的单位元。本章中我们始终假定已作此典范选取，以免追踪大量额外数据。

注意，到目前为止我们只用到 $B$ 上的扇结构，没有用到 $\mathcal{P}$ 的胞腔本身具有格多胞体结构这一事实。另一方面，在第 4 章的情形中，十分重要的是把 $X_0$ 连同对数结构 $X_0^{\dagger}$ 与对数态射 $X_0^{\dagger}\to\operatorname{Spec} k^{\dagger}$（最好是有限光滑的）一并考虑。事实上，格多胞体结构通过决定幽灵层 $\overline{\mathcal{M}}_{X_0}$，给了我们"对数结构应当是什么"的提示。

从第 4 章的例子不难看出这个幽灵层应当是什么。它本质上完全由它在 $X_0$ 各层的一般点处的茎决定。若 $\eta_\tau$ 是 $X_\tau\subseteq X_0$ 的一般点，则在第 4 章由 $(M_{\mathbb{R}},\mathcal{P})$ 构造的退化情形中，有
\begin{equation}
  \mathcal{M}_{X_0,\bar\eta_\tau}=C(\tau)^{\vee}\cap(N\oplus\mathbb{Z}). \tag{6.1}
\end{equation}
事实上，在 $(B,\mathcal{P})$ 的一般情形中我们可以做同样的事情。给定 $\tau\in\mathcal{P}$，$\tau$ 是格多胞体，于是我们可以把它等同于格多胞体 $\tau\subseteq M_{\mathbb{R}}$。此时 $C(\tau)^{\vee}\cap(N\oplus\mathbb{Z})$ 仍然有意义，我们要求 (6.1) 成立。我们还要求 $\overline{\mathcal{M}}_{X_0}$ 在 $X_\tau$ 的大环面轨道上为常值。

结果这一信息足以决定 $\overline{\mathcal{M}}_{X_0}$，但要得到真正的对数结构则困难得多。办法是尝试分类使得
\[
  1\to\mathcal{O}_{X_0}^{\times}\to\mathcal{M}_{X_0}\to\overline{\mathcal{M}}_{X_0}\to0
\]
成为对数结构的扩张，以及给出到 $\operatorname{Spec} k^{\dagger}$ 的态射的截面 $\rho\in\overline{\mathcal{M}}_{X_0}$。这一分类是 [48] 的主要结果。主要困难在于：这样的扩张事实上并非整体存在，而只是在余维二闭子集 $Z\subseteq X_0$ 之外存在。因此 [48] 给出的是在余维二闭子集 $Z\subseteq X_0$ 之外对数光滑的对数态射的描述

% [Tutor Note: (6.1) 的核对锚点：取 $\tau=\{v\}$ 为顶点时，$C(\tau)$ 是 $\Sigma_v$ 的锥，$C(v)^{\vee}\cap(N\oplus\mathbb{Z})$ 正是环面簇 $X_v$ 在对应环面轨道处的对偶幺半群再添一个"$t$ 方向"的坐标；取 $\tau$ 为最大胞腔时则退化为 $N\oplus\mathbb{Z}$。这与第 4 章中心纤维上幽灵层的茎一致。]

% [原文 p.250]
$X_0^{\dagger}\to\operatorname{Spec} k^{\dagger}$。

\begin{example}
坏集 $Z$ 的存在性是理论的基本特征。例如考虑由 $\PP^3\times\mathbb{A}^1$ 中的 $tf_4+x_0\cdot\ldots\cdot x_3=0$ 给出的簇 $X$，其中 $\PP^3$ 的坐标为 $x_0,\dots,x_3$，$\mathbb{A}^1$ 的坐标为 $t$，而 $f_4$ 是 $x_0,\dots,x_3$ 的一般四次齐次多项式。那么投影 $f:X\to\mathbb{A}^1$ 是四次 K3 曲面退化为 $\PP^3$ 中坐标平面之并的退化。令 $X_0=f^{-1}(0)$。

注意 $X$ 恰好在 24 个点
\[
  \{t=f_4=0\}\cap\operatorname{Sing}(X_0)
\]
处奇异。对一般的 $f_4$，这些点是通常二重点。在这些二重点附近 étale 局部地，$X_0\subseteq X$ 看起来像
\[
  V(t)\subseteq\operatorname{Spec} k[x,y,w,t]/(xy-wt).
\]
这个包含在例 3.20 中研究过。因此，若取由除子对数结构 $X_0\subseteq X$ 在 $X_0$ 上诱导的对数结构，则该对数结构在这 24 个点处甚至不是细的。这 24 个点给出坏集 $Z$；显然诱导的映射 $X_0^{\dagger}\to\operatorname{Spec} k^{\dagger}$ 在这里不是对数光滑的。

可以看到，给出这样的 $X_0^{\dagger}$ 的整热带流形由把例 1.28 (3) 的构造应用于自反多胞体
\[
  \Xi=\operatorname{Conv}\{(-1,-1,-1),(1,0,0),(0,1,0),(0,0,1)\}
\]
得到。\hfill$\square$
\end{example}

最后一步是把 $X_0^{\dagger}$ 实现为族 $X\to\operatorname{Spec} k[t]$ 的中心纤维。$X_0^{\dagger}$ 上的对数结构应当是由 $D\cup X_0$ 在 $X$ 上诱导的除子对数结构的拉回，其中 $D\subseteq X$ 是某个除子，其到每个纤维的限制都是反典范除子（于是在 Calabi–Yau 情形 $D$ 为空）。

这实质上就是 [49] 中所解决的问题，其解答将在 §6.2 中就二维情形加以说明。不过，即使尚未解决这个问题，我们仍然可以问：$B$ 中的热带曲线与 $X$ 中的曲线（更具体地说，$X_{\bar\eta}$ 中的曲线，其中 $\eta$ 是 $\operatorname{Spec} k[t]$ 的一般点）之间是否有对应。如第 4 章所提示，为探究这一对应，应先考虑 $B$ 中的热带曲线与 $X_0^{\dagger}$ 中的对数曲线之间的对应。为 $X_0^{\dagger}$ 恰当地定义 Gromov–Witten 不变量，应当随后给出 $X_{\bar\eta}$ 的相同 Gromov–Witten 不变量。这样一来，人们便希望通过热带几何计算 $X_{\bar\eta}$ 的 Gromov–Witten 不变量。

这是目前该纲领中最不成熟的部分，不过我与 Siebert 正在进行的工作 [46]，以及 Parker 的工作 [91]、[92]、[93]，正在发展对数 Gromov–Witten 不变量理论。我们只给出一个非常简单的例子，以展示仿射结构与对数结构的奇点所起的作用。

\begin{example}
我们来考虑例 1.31 (5) 与例 1.33 中讨论的曲面 $(\check B,\check{\mathcal{P}})$。我先断言：这个曲面对应于形如
\[
  X=V(tf_3+x_1x_2x_3)\subseteq\PP^3\times\mathbb{A}^1
\]
的三次曲面退化，其中 $f_3$ 是一般三次齐次多项式。我们有投影 $\pi:X\to\mathbb{A}^1$，

% [原文 p.251]
以及纤维 $X_0=\pi^{-1}(0)$。$X$ 上的对数结构是由 $V(x_0)\cup X_0\subseteq X$ 给出的除子对数结构，它诱导出对数结构 $X_0^{\dagger}$。

为看出 $\check B$ 确实对应于 $X_0^{\dagger}$，令 $X_0'=X_0(\check B,\check{\mathcal{P}})$ 为由 $\check B$ 经本节所述构造得到的簇。先注意 $\check{\mathcal{P}}$ 有三个顶点，而对每个 $v\in\check{\mathcal{P}}$，$\Sigma_v$ 都是 $\PP^2$ 的扇。于是 $\check B$ 给出三个 $\PP^2$ 的副本，它们正是按 $X_0$ 中三个坐标平面相互粘合的方式粘合起来的。所以 $X_0=X_0'$。为确认 $\check B$ 还给出关于 $\overline{\mathcal{M}}_{X_0}$ 的正确信息，我们考察 $X_0$ 各层的一般点。例如，在 $X_0$ 的层 $x=(1:0:0:0)$ 处，$\overline{\mathcal{M}}_{X_0,x}=\mathbb{N}^3$，因为 $V(x_0)\cup X_0$ 在 $x$ 处是正规交叉除子。而 $\check{\mathcal{P}}$ 中相应的极大胞腔是二维标准单形 $\sigma$，且
\[
  C(\sigma)^{\vee}\cap\mathbb{Z}^3\cong\mathbb{N}^3.
\]
类似地，点 $x=(0:1:0:0)$ 是另一个层，它对应于 $\check{\mathcal{P}}$ 中的一个无限矩形 $\sigma$。同样 $\overline{\mathcal{M}}_{X_0,x}=\mathbb{N}^3$，因为 $V(x_0)\cup X_0$ 在 $x$ 处是正规交叉除子，且 $V(x_0)\cup X_0$ 的不可约分支 $V(x_0)$、$V(x_2)$、$V(x_3)$ 都经过 $x$。另一方面，比如写
\[
  \sigma=\operatorname{Conv}\{(0,0),(1,0)\}+\mathbb{R}_{\ge0}(0,1),
\]
则 $C(\sigma)$ 是由 $(0,0,1)$、$(1,0,1)$ 与 $(0,1,0)$ 生成的锥。于是
\[
  C(\sigma)^{\vee}\cap\mathbb{Z}^3\cong\mathbb{N}^3.
\]
类似地可以验证，$\overline{\mathcal{M}}_{X_0}$ 沿一维层与二维层的茎分别与用 $\check{\mathcal{P}}$ 的一维胞腔与零维胞腔所定义的幺半群一致。这说明 $\overline{\mathcal{M}}_{X_0}$ 的茎结构确实是由数据 $(\check B,\check{\mathcal{P}})$ 所指定的那些。还要注意，对一般的 $f_3$，$X$ 的全空间在 $V(f_3)\cap\operatorname{Sing}(X_0)$ 处有九个奇点。这九个奇点构成坏集 $Z$。

再考虑例 1.33 中图示的某条热带曲线。根据第 4 章的经验，我们应当寻找一条映入 $X_0$ 的某个 $\PP^2$ 分支中一条直线的对数曲线，给出 $f:\PP^1\to X_0$。然而，由命题 4.9 我们也知道，当这条 $\PP^1$ 的某个点被映入 $\operatorname{Sing}(X_0)$ 时会有问题：命题 4.9 排除了这种情况。不过，当前情形与第 4 章的情形有一个关键区别：存在轨迹 $Z\subseteq X_0$，在其上对数结构 $X_0^{\dagger}$ 不是细的。$Z$ 的九个点对应于 $\check B$ 中的九个奇点。结果表明，命题 4.9 的论证恰好在如下情形失效：每个满足 $f(p)\in\operatorname{Sing}(X_0)$ 的点 $p\in\PP^1$ 都有 $f(p)\in Z$。

这样的曲线容易计数。在每个平面中，恰有 $3\times3=9$ 条直线满足这一性质。由于有三个平面，这给出 27 种可能选择。其中每一种选择都形变给出 $\pi^{-1}(t)$（$t\neq0$）中的一条真正直线，这正解释了三次曲面上的 27 条直线。在精神上，它们对应于例 1.33 中描述的 27 条热带曲线。这些热带曲线被允许终止于 $\check B$ 的奇点，这反映了如下事实：我们可以有对数曲线 $f:C^{\dagger}\to X_0^{\dagger}$，其不可约分支经过 $Z$ 中的点，而在 $X_0$ 的另一个不可约分支上没有与之匹配的不可约分支。\hfill$\square$

% [Tutor Note: 例 6.1 的 24 个二重点与例 6.2 的三次曲面 27 条直线可作独立核对：$\PP^3$ 中四次 K3 曲面的一般节点型退化给 24 个通常二重点，三次曲面经典地有 27 条直线；两例的坏集 $Z$ 大小（24、9）分别源于这两组数，且 $Z$ 正是对数结构不细之处。]
\end{example}

虽然这一描述非常含糊，但它展示了 A-模型计算的基本策略：$(B,\mathcal{P})$ 给出 $X_0^{\dagger}$，而在

% [原文 p.252]
$(B,\mathcal{P})$ 上作热带曲线计数，应当等价于在 $X_0^{\dagger}$ 上计算对数 Gromov–Witten 不变量，后者又应当等价于在 $X_0^{\dagger}$ 的某个光滑化上的 Gromov–Witten 不变量。不过要注意这一途径的一个特点：若对数 Gromov–Witten 不变量理论被恰当地建立起来，那么 $X_0^{\dagger}$ 上的对数 Gromov–Witten 不变量与光滑化上的 Gromov–Witten 不变量的比较应当是自动的，我们根本不需要知道光滑化的任何信息。

考虑 $X_0$ 上的丰沛线丛 $L$ 也是自然的。特别地，计算 Gromov–Witten 不变量时，我们可能关注 $X_0$ 上给定次数的曲线。这样的线丛的选取给出 $B$ 上的额外数据。事实上，若把 $L$ 限制到 $X_0$ 的任一环面层 $X_\tau$，就得到环面簇 $X_\tau$ 上的丰沛线丛，它由严格凸整分片线性函数 $\varphi_\tau:|\Sigma_\tau|\to\mathbb{R}$ 给出（至多相差一个线性函数）。于是容易验证：若 $\tau_1\subseteq\tau_2$，从而自然地 $X_{\tau_2}\subseteq X_{\tau_1}$，则 $\varphi_{\tau_2}$ 与定义 1.10 中所定义的限制 $\varphi_{\tau_1}(\tau_2)$ 相差一个线性函数。（参见 §3.1.2 中关于线丛限制到环面层的讨论。）

若 $S_\tau:U_\tau\to\mathbb{R}^k$ 定义沿 $\tau$ 的扇结构，则族 $\varphi=\{\varphi_\tau\circ S_\tau\}$ 定义了 $(B,\mathcal{P})$ 上的严格凸多值整分片线性函数。这给出三元组 $(B,\mathcal{P},\varphi)$。

总之，在扇形图像中，数据 $(B,\mathcal{P},\varphi)$ 对应于对 $(X_0^{\dagger},L)$，其中 $X_0^{\dagger}$ 决定 $\mathcal{P}$，而 $L$ 决定 $\varphi$。如前所述，由于粘合的选取，一般会有这样一对的整个模空间，但这里我们忽略这一点。

\subsubsection{锥图像，或 B-模型}

实际上还有另一种解释对 $(B,\mathcal{P})$ 的方式，其灵感来自例 3.6（Mumford 退化）。由一个格多胞体 $\Delta\subseteq N_{\mathbb{R}}$、一个多面体分解 $\mathcal{P}$ 与一个严格凸分片线性函数 $\varphi$ 得到的 Mumford 退化，是 $\PP_{\Delta}$ 的一个退化，其中心纤维是 $\bigcup_{\sigma\in\mathcal{P}_{\max}}\PP_{\sigma}$。

让我们推广这一点，从对 $(B,\mathcal{P})$ 出发。每个 $\sigma\in\mathcal{P}$ 定义一个射影环面簇 $\PP_{\sigma}$。特别地，若 $\sigma_1,\sigma_2\in\mathcal{P}_{\max}$ 且 $\tau=\sigma_1\cap\sigma_2$，我们可以沿公共环面层 $\PP_\tau$ 把 $\PP_{\sigma_1}$ 与 $\PP_{\sigma_2}$ 粘合起来。这样便构造出概形 $\check X_0=\check X_0(B,\mathcal{P})$。同样，如扇形图像中那样，这种粘合有选择的余地，但我们取典范选择，即等同 $\PP_\tau$ 的大环面中的单位元。

每个 $\PP_\sigma$ 还带有丰沛线丛 $\mathcal{O}_{\PP_\sigma}(1)$。同样可以把这些线丛粘合起来（尽管并非对 $\check X_0$ 各分支的每种粘合选择都可行），得到 $\check X_0$ 上的丰沛线丛 $L$。所以本质上数据 $(B,\mathcal{P})$ 决定了对 $(\check X_0,L)$。然而它并不决定对数结构；后者需要额外选取 $B$ 上的多值严格凸分片线性函数 $\varphi$，这与 Mumford 退化相类比。

给定这样的函数，它如下决定幺半群 $\mathcal{M}_{\check X_0,\bar\eta_\tau}$，其中 $\eta_\tau$ 是 $\PP_\tau$（$\tau\in\mathcal{P}$）的一般点。至多相差仿射线性函数，$\varphi$ 在 $\tau$ 的邻域内由分片线性函数 $\varphi_\tau:|\Sigma_\tau|\to\mathbb{R}$ 决定，其中 $\Sigma_\tau$ 是由沿 $\tau$ 的扇结构决定的 $\mathbb{R}^k$ 中的扇，$k$ 是 $\tau$ 的余维。于是
\begin{equation}
  \mathcal{M}_{\check X_0,\bar\eta_\tau}=\{(m,r)\mid m\in|\Sigma_\tau|,\ r\ge\varphi_\tau(m)\}\cap(\mathbb{Z}^k\oplus\mathbb{Z}). \tag{6.2}
\end{equation}

% [原文 p.253]
如同扇形图像中那样，这些茎决定了 $\overline{\mathcal{M}}_{\check X_0}$，然后按照同样的程序分类由 $\mathcal{O}_{\check X_0}^{\times}$ 扩张 $\overline{\mathcal{M}}_{\check X_0}$ 得到的、给出对数结构的 $\mathcal{M}_{\check X_0}$。

于是我们看到了扇形图像与锥图像之间的对称性：在扇形图像中，$\mathcal{P}$ 决定 $M$、$\varphi$ 决定 $L$；而在锥图像中，$\mathcal{P}$ 决定 $L$、$\varphi$ 决定 $M$。

现在在 Calabi–Yau 情形，B-模型涉及度量复结构变分的周期积分。要谈论这些，我们需要以 $\check X_0^{\dagger}$ 为中心纤维的族 $\check X\to\operatorname{Spec} k[t]$。因此仅仅得到 $\check X_0^{\dagger}$ 是不够的，我们需要整个族。这就是我们所需的工作背景：

\begin{question}
给定给出 $\check X_0^{\dagger}$ 的 $(B,\mathcal{P},\varphi)$，是否存在以 $\check X_0$ 为中心纤维的族 $\check X\to\operatorname{Spec} k[t]$，以及反典范除子的族 $D\subseteq\check X$，它诱导除子对数结构 $D\cup\check X_0\subseteq\check X$，并在中心纤维上限制为对数结构 $\check X_0^{\dagger}$？
\end{question}

[49] 的主要结果给出了 $\check X_0^{\dagger}$ 使该问题有肯定解答的充分条件。此外，我们将看到，这个解答具有非常热带的性质。

在 §6.2 中，我们将给出 $\dim B=2$ 时构造族 $\check X\to\operatorname{Spec} k[t]$ 的论证的大部分细节；高维情形要困难得多。

% [Tutor Note: §6.1 的两个图像互为"对偶"：扇形图像中 $\mathcal{P}\mapsto M$、$\varphi\mapsto L$；锥图像中 $\mathcal{P}\mapsto L$、$\varphi\mapsto M$。核对点：(6.1) 用锥 $C(\tau)^{\vee}$，(6.2) 用 $\varphi_\tau$ 的 Newton 多面体上方区域，二者在离散 Legendre 变换下互换——这正是 §6.1.3 的论证。]

\subsubsection{镜像对称与离散 Legendre 变换}

粗略地说，前两节论证了：给数据 $(B,\mathcal{P},\varphi)$ 可以配上由扇形图像的构造得到对数空间 $X_0(B,\mathcal{P})^{\dagger}$，以及由锥图像的构造得到 $\check X_0(B,\mathcal{P})^{\dagger}$。如我们所见，粘合与对数结构都有选择余地，所以实际上存在这类结构的模空间；但为简单起见，这里我们忽略它。

上述讨论提示了如下镜像对称的程序：

(1) 从"良好"的退化 $\pi:X\to\operatorname{Spec} k[t]$ 出发，即中心纤维是环面簇之并、且 $\pi$ 在某个集合 $Z$ 之外对数光滑。该集合应当是余维二的，且不含任何环面层。还有一些更技术性的条件。这里我们不给出精确的条件集，完整定义见 [48]，定义 4.1 或 [41]，定义 7.1。简言之，我们希望 $\pi$ 是我们所谓的\emph{环面退化}。此外，选取 $X$ 上相对丰沛的线丛 $L$。

(2) $X_0=\pi^{-1}(0)$ 上的对数结构决定一个带奇点的整热带流形 $(B,\mathcal{P})$，使得 $X_0=X_0(B,\mathcal{P})$（在扇形图像中）；而 $L$ 决定一个多值整分片线性严格凸函数 $\varphi$，给出三元组 $(B,\mathcal{P},\varphi)$。

(3) 现在在锥图像中看待 $(B,\mathcal{P},\varphi)$，它决定了带对数结构的 $\check X_0(B,\mathcal{P})$。构造以 $\check X_0(B,\mathcal{P})$ 为中心纤维的族 $\check X\to\operatorname{Spec} k[t]$。族 $X,\check X\to\operatorname{Spec} k[t]$ 应当是互为镜像的。

有时我们可能希望只在一侧考虑 A-模型与 B-模型，比如对 $X\to\operatorname{Spec} k[t]$。除了 $X_0=X_0(B,\mathcal{P})$、$L$ 决定 $\varphi$ 的三元组 $(B,\mathcal{P},\varphi)$ 之外，还有另一个三元组 $(\check B,\check{\mathcal{P}},\check\varphi)$，使得 $X_0^{\dagger}=\check X_0(\check B,\check{\mathcal{P}})^{\dagger}$。

% [原文 p.254]
$(B,\mathcal{P},\varphi)$ 与 $(\check B,\check{\mathcal{P}},\check\varphi)$ 之间是什么关系？记号本身已经透露了答案：它们由 §1.5 的离散 Legendre 变换相联系。确实，$\mathcal{P}$ 与 $\check{\mathcal{P}}$ 之间存在反序对应，记作 $\tau\mapsto\check\tau$，因为 $\mathcal{P}$ 与 $\check{\mathcal{P}}$ 都与 $X_0$ 的各层一一对应，即 $X_\tau=\PP_{\check\tau}$。扇 $\Sigma_\tau$ 正是 $\check\tau$ 的法扇，这恰是因为 $X_\tau=\PP_{\check\tau}$。此外，$\check\tau$ 是 $\varphi_\tau$ 的 Newton 多面体。又 $\Sigma_{\check\tau}$ 承载函数 $\check\varphi_{\check\tau}$，它通过 (6.2) 决定 $\check X_0(\check B,\check{\mathcal{P}})$ 上的对数结构。比较 (6.2) 与 (6.1)，我们看到需要两个锥 $C(\tau)^{\vee}$ 与 $\{(m,r)\mid m\in|\Sigma_{\check\tau}|,\ r\ge\check\varphi_{\check\tau}(m)\}$ 同构。容易验证这成立当且仅当 $\tau$ 是 $\check\varphi_{\check\tau}$ 的 Newton 多面体（的平移）。与 §1.5 的构造相比较，便知 $(\check B,\check{\mathcal{P}},\check\varphi)$ 确实是 $(B,\mathcal{P},\varphi)$ 的离散 Legendre 变换。

所以，在某一层面上，离散 Legendre 变换实现了镜像对称。

\begin{example}
让我们从这个观点来看 $\PP^2$ 的镜像。$\PP^2\times\mathbb{A}^1_k\to\mathbb{A}^1_k$ 是 $\PP^2$ 退化为单个环面簇的（平凡！）退化。$(B,\mathcal{P})$ 就是 $(\mathbb{R}^2,\Sigma_{\PP^2})$，其中 $\Sigma_{\PP^2}$ 是 $\PP^2$ 的扇。选取 $\PP^2$ 上的丰沛线丛 $\mathcal{O}_{\PP^2}(n)$ 决定（至多相差线性函数）一个整分片线性函数 $\varphi:B\to\mathbb{R}$，其 Newton 多面体（至多相差平移）是 $\operatorname{Conv}\{(0,0),(n,0),(0,n)\}$。这个 Newton 多面体就是 $\check B$，而 $\check{\mathcal{P}}$ 是 $\check B$ 的面，$\check\varphi$ 是线性的。取 $n=1$。

在锥图像中看待 $(B,\mathcal{P},\varphi)$，我们可以直接构造 Mumford 退化。这里
\[
  \tilde\Delta=\{(m,r)\mid m\in\mathbb{R}^2,\ r\ge\varphi(m)\},
\]
并且
\begin{align*}
  C(\tilde\Delta)&=\{(sm,sr,s)\mid s\ge0,\ (m,r)\in\tilde\Delta\}\\
  &=\{(m,r,s)\mid s\ge0,\ (m,r)\in\tilde\Delta\},
\end{align*}
因为 $\tilde\Delta$ 在伸缩下不变。于是
\[
  C(\tilde\Delta)\cap(\mathbb{Z}^2\oplus\mathbb{Z}\oplus\mathbb{Z})=(\tilde\Delta\cap(\mathbb{Z}^2\oplus\mathbb{Z}))\oplus\mathbb{N},
\]
且
\[
  \operatorname{Proj}k[(\tilde\Delta\cap(\mathbb{Z}^2\oplus\mathbb{Z}))\oplus\mathbb{N}]
  \cong\operatorname{Spec}k[\tilde\Delta\cap(\mathbb{Z}^2\oplus\mathbb{Z})].
\]
于是令
\[
  P=\{(m,r)\mid m\in\mathbb{Z}^2,\ r\ge\varphi(m)\}\subseteq\mathbb{Z}^2\oplus\mathbb{Z},
\]
我们得到自然映射
\[
  \operatorname{Spec} k[P]\to\operatorname{Spec} k[\mathbb{N}]=\mathbb{A}^1_k,
\]
它由自然包含 $\mathbb{N}\hookrightarrow P$（$n\mapsto(0,0,n)$）诱导。注意 $P\cong\mathbb{N}^3$：取恰当的 $\varphi$，我们可以取 $P$ 的生成元为 $(1,0,0)$、$(0,1,0)$ 与 $(-1,-1,1)$，而映射可写成
\[
  \operatorname{Spec} k[x_0,x_1,x_2]\to\operatorname{Spec} k[\kappa],\qquad
  \kappa\mapsto x_1x_1x_2.
\]
这给出 §2.2 中 $\PP^2$ 的镜像族。这一做法对其他环面簇也有效，但要得到 §5.1.1 中的镜像构造，需要这一构造的一个变体来获得多维族。\hfill$\square$
\end{example}

% [Tutor Note: 例 6.4 末的 $\kappa\mapsto x_1x_1x_2$ 是原文印刷笔误：按生成元 $(1,0,0),(0,1,0),(-1,-1,1)$ 与 $n\mapsto(0,0,n)$ 计算，$t$ 对应三个生成元之和，故应为 $\kappa\mapsto x_0x_1x_2$。本译照原文抄录并在此注明。]

% [原文 p.255]
\subsection{从整热带流形到二维退化}

本节固定三元组 $(B,\mathcal{P},\varphi)$，其中
\begin{itemize}
\item $(B,\mathcal{P})$ 是二维整热带流形，可能带边界，也可能非紧；
\item $\varphi$ 是带整斜率的分片线性多值严格凸函数。
\end{itemize}
关于 $B$ 的奇点，我们还需作两条进一步假设，它们涉及局部系统 $\Lambda$ 绕每个奇点的单值。回忆 $\Lambda$ 是由 $B_0$ 上的整仿射结构决定的、$B_0\subseteq B$ 上整向量场的局部系统。

由构造，这些奇点只出现在 $\mathcal{P}$ 的不含于 $B$ 的边界的紧边上。设 $\omega\in\mathcal{P}$ 是这样一条边，它含于两个二维胞腔 $\sigma_{\pm}\in\mathcal{P}$ 中。把 $\omega$ 的顶点标为 $v_+$ 与 $v_-$。设 $d_\omega\in\Lambda_{v_+}$ 是 $\omega$ 在 $v_+$ 处的本原切向量，指向 $v_-$。设 $\check d_\omega\in\check\Lambda_{v_+}$ 是唯一的整余切向量，它本原、零化 $d_\omega$，且在指向 $\sigma_+$ 内部的切向量上取正值。我们也可以把 $d_\omega$ 看作在指定 $\omega$ 的哪个顶点是 $v_+$、哪个是 $v_-$；类似地，$\check d_\omega$ 指定哪个胞腔是 $\sigma_+$、哪个是 $\sigma_-$。

现在考虑以 $v_+$ 为基点、依次经过 $\sigma_+$、$v_-$、$\sigma_-$ 再回到 $v_+$ 的环路 $\gamma$，如图 1 所示。那么不难看出，$\Lambda$ 绕 $\gamma$ 的单值——它是线性变换
\[
  T_\gamma:\Lambda_{v_+}\to\Lambda_{v_+},
\]
——具有形式
\[
  T_\gamma(m)=m+\kappa_\omega\langle\check d_\omega,m\rangle d_\omega,
\]
其中 $\kappa_\omega$ 是某个整数。（见例 1.28 (4)，那里就一个显式情形算出了这一点。）注意 $\kappa_\omega$ 与选择无关：交换 $v_+$ 与 $v_-$ 或 $\sigma_+$ 与 $\sigma_-$ 会改变 $d_\omega$ 或 $\check d_\omega$ 的符号，但同时也会反转 $\gamma$ 的方向。

\begin{definition}
若对 $\mathcal{P}$ 中所有紧且不含于 $\partial B$ 的边 $\omega$ 都有 $\kappa_\omega\ge0$，则称 $B$ 是\emph{正}的。若对所有这样的边 $\omega$ 都有 $\kappa_\omega\in\{0,1\}$，则称 $B$ 是\emph{简单}的。
\end{definition}

\begin{figure}[htbp]
\centering
\includegraphics[width=0.42\textwidth]{f6_1.png}
\figtitle{1}{两个二维胞腔 $\sigma_+,\sigma_-$ 沿公共边 $v_+v_-$ 粘合，虚线为环路 $\gamma$。}
\end{figure}

% [原文 p.256]
正性即使对 $\check X_0(B,\mathcal{P})$ 上对数结构的存在性也是必需的；若正性不成立，对数结构会在 $\check X_0(B,\mathcal{P})$ 的整个层上变坏，那时光滑化便毫无指望。简单性并非必需，但它会使我们的论证简单一些。

以下我们假定 $(B,\mathcal{P})$ 是正且简单的。现在的目标是构造 $\check X_0(B,\mathcal{P})$ 的一个显式光滑化。特别地，我们将证明如下形式陈述：

\begin{theorem}
设 $(B,\mathcal{P})$ 正且简单，$\dim B=2$，并给定 $(B,\mathcal{P})$ 上带整斜率的多值严格凸整分片线性函数 $\varphi$，则可以构造形式平坦光滑化 $\check X(B,\mathcal{P},\varphi)\to\operatorname{Spf} k[t]$。
\end{theorem}

所谓形式，是指 $\check X(B,\mathcal{P},\varphi)$ 是形式概形，其底拓扑空间是 $\check X_0(B,\mathcal{P})$。形式概形 $\operatorname{Spf} k[t]$ 是仿射直线在原点处的完备化。

为得到 $\operatorname{Spec} k[t]$ 上真正的平坦形变 $\check X(B,\mathcal{P},\varphi)$，我们还需要一个要素。只要 $\check X(B,\mathcal{P},\varphi)\to\operatorname{Spf} k[t]$ 是真态射（当 $B$ 紧时如此），且 $\check X(B,\mathcal{P},\varphi)$ 上存在相对丰沛线丛 $L$，Grothendieck 存在定理（[52]，5.4.5）就告诉我们：$\check X(B,\mathcal{P},\varphi)$ 是某个 $\operatorname{Spec} k[t]$ 上的概形 $\check X(B,\mathcal{P},\varphi)$ 沿中心纤维（同构于 $\check X_0(B,\mathcal{P})$）的形式完备化。平坦族 $\check X(B,\mathcal{P},\varphi)\to\operatorname{Spec} k[t]$ 就是所要的族。

% [Tutor Note: 定理 6.6 只给出"形式"平坦光滑化；要得到 $\operatorname{Spec} k[t]$ 上真正的族，需 Grothendieck 存在定理（[52]，5.4.5）与相对丰沛线丛，正文随即说明。"正的"条件 $\kappa_\omega\ge0$ 是 $\check X_0(B,\mathcal{P})$ 上对数结构存在的必要条件；"简单"（$\kappa_\omega\in\{0,1\}$）只是使论证简化。]

如先前所述，$\check X_0(B,\mathcal{P})$ 带有自然的丰沛线丛，它由把各不可约分支 $\PP_\sigma$（$\sigma\in\mathcal{P}_{\max}$）上的线丛 $\mathcal{O}_{\PP_\sigma}(1)$ 粘合而成。事实上可以验证它延拓为 $\check X(B,\mathcal{P},\varphi)$ 上的丰沛线丛；这只是一个技术要点，这里不再赘述。我们将集中证明定理 6.6，它是 [49] 主要结果的特例。

\subsubsection{预热：Mumford 退化}

为说明基本策略，我们先更详细地研究 Mumford 退化的情形。从 $B\subseteq N_{\mathbb{R}}$（格多面体）、把 $B$ 分解为格多胞体的多面体分解 $\mathcal{P}$，以及带整斜率的凸分片线性函数 $\varphi:B\to\mathbb{R}$ 出发。由例 3.6 可知，这决定一个多面体
\[
  \tilde\Delta=\{(n,r)\in N_{\mathbb{R}}\oplus\mathbb{R}\mid n\in B,\ r\ge\varphi(n)\},
\]
它又定义环面簇 $\PP_{\tilde\Delta}$ 以及映射
\[
  \pi:\PP_{\tilde\Delta}\to\mathbb{A}^1_k.
\]
这个途径不能立刻对更有趣的 $B$ 生效，其基本原因是这一构造太整体了，特别地它要求把 $B$ 嵌入 $N_{\mathbb{R}}$。一般地，我们只有 $B$ 在奇点之外的局部嵌入。因此，在推广这一想法之前，我们首先需要把构造局部化。

第一步是考虑它的无穷小版本。令
\[
  \mathcal{O}_k:=\operatorname{Spec} k[t]/(t^{k+1}),
\]
考虑
\[
  \pi_k:X_k:=\PP_{\tilde\Delta}\times_{\mathbb{A}^1}\mathcal{O}_k\to\mathcal{O}_k.
\]

% [原文 p.257]
这样做之所以合理，一个原因是：一般地我们只要求 $\operatorname{Spec} k[t]$ 上的退化，而非 $\mathbb{A}^1_k$ 上的退化，且一般也不指望更好的结果（见例 6.15）。于是我们可以逐阶构造这个形变，再取极限。这正是我们在一般情形中要做的。

下一步是详细描述 $X_k$ 的仿射覆盖。设 $\check\Sigma$ 是 $\tilde\Delta$ 的法扇，故 $\PP_{\tilde\Delta}=X_{\check\Sigma}$。当然 $X_{\check\Sigma}$ 有以 $\check\Sigma$ 的锥为指标的仿射覆盖。特别地，若 $\tau\in\mathcal{P}$ 且
\[
  \tilde\tau:=\{(n,r)\mid n\in\tau,\ r=\varphi(n)\}\subseteq\tilde\Delta
\]
是 $\tilde\Delta$ 中投影到 $\tau$ 的相应水平面，则用 $\check\tau$ 表示法锥 $N_{\tilde\Delta}(\tilde\tau)$。注意 $\check\tau$ 决定 $X_{\check\Sigma}$ 的一个仿射开子集，我们记作 $U_\tau$。则 $\{U_\tau\mid\tau\in\mathcal{P}\}$ 构成 $\PP_{\tilde\Delta}$ 的开覆盖。确实，当 $\tau$ 遍历 $\mathcal{P}$ 的胞腔时 $\check\tau$ 并不遍历 $\check\Sigma$ 的所有锥，但 $\check\Sigma$ 的所有极大锥都会出现，因为极大锥对偶于 $\tilde\Delta$ 的零维面；而由扇 $\Sigma$ 的极大锥对应的 $X_\Sigma$ 仿射子集覆盖 $X_\Sigma$，便知 $\{U_\tau\mid\tau\in\mathcal{P}\}$ 覆盖 $\PP_{\tilde\Delta}$。

开集 $U_\tau$ 可以如下描述，利用事实 (1.3)：
\[
  (\check\tau)^{\vee}=(N_{\tilde\Delta}(\tilde\tau))^{\vee}=T_{\tilde\tau}\tilde\Delta.
\]
（另见注 3.5。）我们可以如下描述 $T_{\tilde\tau}\tilde\Delta$。对 $\sigma\in\mathcal{P}_{\max}$，设 $m_\sigma\in M$ 是 $\varphi|_\sigma$ 的斜率，即
\[
  m_\sigma=d(\varphi|_\sigma)\in M.
\]
定义 $\varphi_\tau:N_{\mathbb{R}}\to\mathbb{R}$ 为
\[
  \varphi_\tau(n):=\max\{\langle m_\sigma,n\rangle\mid\tau\subseteq\sigma\in\mathcal{P}_{\max}\}.
\]
则容易验证
\[
  T_{\tilde\tau}\tilde\Delta=\{(n,r)\in N_{\mathbb{R}}\oplus\mathbb{R}\mid r\ge\varphi_\tau(n)\}.
\]
令
\[
  P_\tau:=T_{\tilde\tau}\tilde\Delta\cap(N\oplus\mathbb{Z})
  =\{(n,r)\in N\oplus\mathbb{Z}\mid r\ge\varphi_\tau(n)\}.
\]
则由定义
\[
  U_\tau=\operatorname{Spec} k[P_\tau].
\]
当然 $\pi$ 限制为映射 $\pi:U_\tau\to\mathbb{A}^1_k$，它由正则函数 $t=z^{(0,1)}$ 给出；然后我们作基变换
\[
  \pi_k:U_\tau^k=U_\tau\times_{\mathbb{A}^1}\mathcal{O}_k\to\mathcal{O}_k.
\]
描述集合 $U_\tau^k$ 现在是 $B$ 上的一个更局部的问题。不过，如我们将看到的，它还不够局部。为进一步细化这一描述，注意 $U_\tau^k$ 与 $U_\tau^0=\pi^{-1}(0)\cap U_\tau$ 有相同的底拓扑空间，从而有相同的不可约分支。我们希望把 $U_\tau^k$ 的各不可约分支粘合起来描述它。具体做法如下。我们首先作定义

% [原文 p.258]
\begin{definition}
设 $p=(n,r)\in P_\tau$，$\tau\subseteq\sigma\in\mathcal{P}_{\max}$。则 $p$ 沿 $\sigma$ 的\emph{阶}为
\[
  \operatorname{ord}_\sigma(p)=r-\langle m_\sigma,n\rangle.
\]
这恰是单项式 $z^p\in k[P_\tau]$ 沿 $U_\tau$ 中由 $\tilde\sigma$ 指定的环面除子的消没阶。

用 $\mathcal{P}^{\partial}_{\max}$ 表示 $\mathcal{P}$ 中维数为 $\dim B-1$ 且含于 $\partial B$ 的胞腔之集。对 $\tau\subseteq\sigma\in\mathcal{P}^{\partial}_{\max}$，设 $m_\sigma\in M$ 是射线 $N_{\Delta}(\sigma)$ 的本原生成元；注意它既是 $\Delta$ 的法扇中的射线，也是 $\tilde\Delta$ 的法扇中的射线。则定义
\[
  \operatorname{ord}_\sigma(p)=\langle m_\sigma,n\rangle.
\]
这恰是单项式 $z^p\in k[P_\tau]$ 沿 $\tilde\Delta$ 的极大面 $\{(n,r)\in\tilde\Delta\mid n\in\sigma\}$ 所指定的环面除子的消没阶。

给定对 $\omega\subseteq\tau$（$\omega,\tau\in\mathcal{P}$），令
\[
  R_\omega:=k[P_\omega].
\]
考虑由
\[
  I^k_{\omega,\tau}:=\{p\in P_\omega\mid\exists\sigma\in\mathcal{P}_{\max}\cup\mathcal{P}^{\partial}_{\max},\ \tau\subseteq\sigma,\ \operatorname{ord}_\sigma(p)>k\}
\]
定义的单项式理想 $I^k_{\omega,\tau}\subseteq P_\omega$。$I^k_{\omega,\tau}$ 于是生成 $R_\omega$ 中的理想，我们仍记作 $I^k_{\omega,\tau}$，并定义
\[
  R^k_{\omega,\tau}:=R_\omega/I^k_{\omega,\tau}.
\]
\hfill$\square$
\end{definition}

% [Tutor Note: $\operatorname{ord}_\sigma$ 的两种定义要与 $t=z^{(0,1)}$ 的约定配套：内部胞腔用斜率 $m_\sigma=d(\varphi|_\sigma)$，边界极大胞腔用 $N_\Delta(\sigma)$ 的本原生成元（不含 $t$ 方向），故后者的阶只取 $\langle m_\sigma,n\rangle$。$I^k_{\omega,\tau}$ 的条件是 $>k$ 而非 $\ge k$，对应"至多 $k$ 阶加厚" $R^k_{\omega,\tau}$。]

注意由 $\varphi_\omega$ 的凸性，
\[
  I^0_{\omega,\tau}=P_\omega\setminus\{(n,r)\mid n\in T_\omega\tau,\ r=\varphi_\omega(n)\}.
\]
于是
\[
  R^0_{\omega,\tau}\cong k[(T_\omega\tau)\cap N].
\]
由于 $T_\omega\tau=(N_\tau(\omega))^{\vee}$（由 (1.3)）且 $N_\tau(\omega)$ 是 $\tau$ 的法扇 $\check\Sigma_\tau$ 中的锥，事实上可以把 $\operatorname{Spec} R^0_{\omega,\tau}$ 看作 $\PP_\tau$（$\pi^{-1}(0)$ 的环面层）中由面 $\omega\subseteq\tau$ 决定的仿射子集。于是可以把 $\operatorname{Spec} R^k_{\omega,\tau}$ 看作 $U_\omega$ 中闭子集 $\operatorname{Spec} R^0_{\omega,\tau}$ 的一种加厚。

给定 $\omega\subseteq\tau_1\subseteq\tau_2$（$\omega,\tau_1,\tau_2\in\mathcal{P}$），有包含 $I^k_{\omega,\tau_2}\subseteq I^k_{\omega,\tau_1}$，从而有满射
\[
  \psi_{\tau_1,\tau_2}:R^k_{\omega,\tau_2}\twoheadrightarrow R^k_{\omega,\tau_1}.
\]

\begin{example}
在图 2 中我们画出一个一维例子。这里 $B=[0,4]\subseteq N_{\mathbb{R}}=\mathbb{R}$，而 $\varphi$ 的图像是图中 $\tilde\Delta$ 的下部。右侧画出 $\tilde\Delta$ 的法扇 $\check\Sigma$。图 3 展示了不同 $\omega$ 下的 $\varphi_\omega$ 与 $P_\omega$。图 4 展示了不同选择下的理想 $I^k_{\omega,\tau}$。\hfill$\square$
\end{example}

% [原文 p.259]
\begin{figure}[htbp]
\centering
\includegraphics[width=0.78\textwidth]{f6_2.png}
\figtitle{2}{一维例子：$\varphi$ 的图像（左）与 $\tilde\Delta$ 的法扇 $\check\Sigma$（右）。}
\end{figure}

固定 $\omega$，环的集合
\[
  \{R^k_{\omega,\tau}\mid\omega\subseteq\tau\},
\]
连同映射 $\psi_{\tau_1,\tau_2}$，构成一个逆向系统。我们有

\begin{lemma}
\[
  \lim_{\longleftarrow}R^k_{\omega,\tau}\cong R_\omega/(t^{k+1}),
\]
其中 $t=z^{\rho}$，$\rho=(0,1)\in P_\omega$。
\end{lemma}

% [Tutor Note: Lemma 6.9 的证明跨页到 p.262（下一块 c06b 开头 "To show this independence …"），故本块此处只给出证明的开头、未闭合 proof 环境；c06b 续写后在论证末尾用 \hfill$\square$ 收尾即可。]

%（p.259 同页续）
\para{证明。}逆向极限中的一个元素是一族元素 $(f_\tau)_{\omega\subseteq\tau}$，其中 $f_\tau\in R^k_{\omega,\tau}$，且只要 $\tau_1\subseteq\tau_2$，就有
\[
  \psi_{\tau_1,\tau_2}(f_{\tau_2})=f_{\tau_1}.
\]
由于对所有 $\tau\supseteq\omega$ 有 $t^{k+1}\in I^k_{\omega,\tau}$，我们有自然满射
\[
  \psi_\tau:R_\omega/(t^{k+1})\to R^k_{\omega,\tau}.
\]
于是显然映射
\[
  f\mapsto(\psi_\tau(f))_{\omega\subseteq\tau}
\]
定义同态
\[
  \psi:R_\omega/(t^{k+1})\to\lim_{\longleftarrow}R^k_{\omega,\tau},
\]
我们只需证明它是同构。

对单射性，设 $\psi(\sum_p c_pz^p)=0$，其中 $p$ 遍历 $P_\omega\setminus((k+1)\rho+P_\omega)$ 的元素。这只可能在只要 $c_p\neq0$ 就有 $\psi(z^p)=0$ 时发生。但若 $\psi(z^p)=0$，则对所有 $\tau\supseteq\omega$ 有 $p\in I^k_{\omega,\tau}$。特别地，取 $\tau=\sigma\in\mathcal{P}_{\max}$，若写 $p=(n,r)$，就必须有 $r>\langle m_\sigma,n\rangle+k$。于是对所有含 $\omega$ 的 $\sigma\in\mathcal{P}_{\max}$，有 $r-(k+1)\ge\langle m_\sigma,n\rangle$，从而 $r-(k+1)\ge\varphi_\omega(n)$。因此 $p-(k+1)\rho\in P_\omega$，即 $z^p\in(t^{k+1})$。这证明单射性。

对满射性，考虑逆向极限中的 $(f_\tau)_{\omega\subseteq\tau}$。于是可写
\[
  f_\tau=\sum_{p\in P_\omega\setminus I^k_{\omega,\tau}}c_{p,\tau}z^p.
\]
设 $p\in P_\omega\setminus I^k_{\omega,\tau}$ 且 $p\in P_\omega\setminus I^k_{\omega,\tau'}$，其中 $\tau,\tau'\supseteq\omega$ 不同。若我们能证明 $c_{p,\tau}=c_{p,\tau'}$，从而可定义与 $\tau$ 无关的 $c_p=c_{p,\tau}$，便可由
\[
  f=\sum_{p\in P_\omega\setminus\bigcap_\tau I^k_{\omega,\tau}}c_pz^p
\]
定义 $f\in R_\omega/(t^{k+1})$，使得 $\psi(f)=(f_\tau)_{\omega\subseteq\tau}$，从而说明满射性。

% [原文 p.260]
\begin{figure}[htbp]
\centering
\includegraphics[width=0.34\textwidth]{f6_3.png}
\figtitle{3}{对 $\omega=\{1\},\{2\},[2,4],\{0\}$ 分别画出的 $\varphi_\omega$ 与 $P_\omega$。}
\end{figure}

% [原文 p.261]
\begin{figure}[htbp]
\centering
\includegraphics[width=0.34\textwidth]{f6_4.png}
\figtitle{4}{阴影区域表示理想中幺半群的点。}
\end{figure}

% [Tutor Note: 块至此为 p.261（Figure 4 所在页）；Lemma 6.9 的证明余下部分从 p.262 起（"To show this independence…"），由 c06b 续译并以 □ 收尾。图中标签 $I^2_{\{1\},[0,1]}$ 等沿用原文记号，未重画。]
